\documentclass[12pt,a4paper,oneside]{book}
\usepackage[utf8]{inputenc}
\usepackage[T1]{fontenc}
\usepackage[english,french]{babel}
\usepackage{graphicx}
\usepackage{geometry}
\usepackage[table]{xcolor}
\usepackage{colortbl}
\usepackage{titlesec}
\usepackage{fancyhdr}
\usepackage{hyperref}
\usepackage{xcolor}
\usepackage{listings}
\usepackage{multicol}
\usepackage{amsmath}
\usepackage{amssymb}
\usepackage{booktabs}
\usepackage{longtable}
\usepackage{parskip}
\usepackage{enumitem}
\usepackage{fontawesome5}
\usepackage{url}

\definecolor{noir}{RGB}{9, 9, 11}
\definecolor{bleuIA}{RGB}{37, 99, 235}
\definecolor{violet}{RGB}{124, 58, 237}
\definecolor{cyan}{RGB}{6, 182, 212}
\definecolor{or}{RGB}{212, 175, 55}
\definecolor{vert}{RGB}{34, 197, 94}
\definecolor{gris}{RGB}{107, 114, 128}

\geometry{top=2.5cm,bottom=2.5cm,left=2.5cm,right=2.5cm}

\hypersetup{
    colorlinks=true,
    linkcolor=bleuIA,
    filecolor=violet,
    urlcolor=cyan,
    pdfauthor={E-Graphisme by ELECTRON},
    pdftitle={Documentation Enterprise - Volume 9},
    pdfsubject={API Enterprise Platform}
}

\pagestyle{fancy}
\fancyhf{}
\fancyhead[LE,RO]{\thepage}
\fancyhead[RE,LO]{\leftmark}
\fancyfoot[CE,CO]{\thepage}

\titleformat{\chapter}[display]
{\normalfont\Huge\bfseries\color{bleuIA}}
{\chaptertitlename\ \thechapter}{20pt}{\Huge}
\titleformat{\section}
{\normalfont\Large\bfseries\color{bleuIA}}
{\thesection}{1em}{}
\titleformat{\subsection}
{\normalfont\large\bfseries\color{violet}}
{\thesubsection}{1em}{}

\begin{document}

\begin{titlepage}
\thispagestyle{empty}
\begin{center}
\vspace*{2cm}
{\Huge\bfseries\color{bleuIA}E-GRAPHISME}\\[0.3cm]
{\LARGE\bfseries\color{violet}by ELECTRON}\\[2cm]
\vspace{2cm}
{\huge\bfseries\color{noir}API ENTERPRISE PLATFORM}\\[1cm]
{\Large\bfseries\color{cyan}Volume 9}\\[3cm]
\vspace{2cm}
{\Large\color{gris}Une plateforme ouverte, sécurisée et entièrement interconnectée}\\[3cm]
{\Large\bfseries\color{or}Version 2026}\\[2cm]
\vfill
\end{center}
\end{titlepage}

\frontmatter

\chapter*{Préface}
\addcontentsline{toc}{chapter}{Préface}
Ce volume définit l'architecture complète des API et des intégrations de E-Graphisme by ELECTRON.

\bigskip
\textbf{Classification :} Document Interne - Partageable\\
\textbf{Volume :} 9\\
\textbf{Version :} 2026

\tableofcontents

\mainmatter

\chapter{VISION DE L'ARCHITECTURE}

\section{Principes}
Toutes les communications doivent être centralisées. Aucun module ne dialogue directement avec un autre.

\section{Architecture}
\begin{enumerate}
\item Frontend
\item API Gateway
\item Microservices
\item Database
\item AI Agents
\item External Services
\end{enumerate}

\chapter{ARCHITECTURE MICROSERVICES}

\section{Services}
\begin{multicols}{2}
\begin{itemize}
\item Auth Service
\item User Service
\item Company Service
\item CRM Service
\item ERP Service
\item Marketplace Service
\item Boutique Service
\item AI Service
\item Media Service
\item Billing Service
\item Affiliate Service
\item Analytics Service
\item Notification Service
\item Delivery Service
\end{itemize}
\end{multicols}

\chapter{API GATEWAY}

\section{Fonctions}
\begin{multicols}{2}
\begin{itemize}
\item Routage
\item Authentification
\item Limitation du trafic
\item Journalisation
\item Compression
\item Cache
\item Versioning
\item Surveillance
\item Gestion des erreurs
\end{itemize}
\end{multicols}

\chapter{SERVICE REGISTRY}

\section{Fonctions}
\begin{itemize}
\item Enregistrement automatique
\item Découverte des services
\item Vérification disponibilité
\item Surveillance état
\end{itemize}

\chapter{AUTHENTICATION API}

\section{Fonctions}
\begin{multicols}{2}
\begin{itemize}
\item Connexion
\item Déconnexion
\item Création de compte
\item Vérification email
\item Réinitialisation mot de passe
\item Double authentification (2FA)
\item Gestion des sessions
\end{itemize}
\end{multicols}

\chapter{AUTHORIZATION API}

\section{Gestion}
\begin{multicols}{2}
\begin{itemize}
\item RBAC
\item Permissions
\item Rôles
\item Groupes
\item Organisations
\item Droits des agents IA
\end{itemize}
\end{multicols}

\chapter{USERS API}

\section{Opérations}
\begin{multicols}{2}
\begin{itemize}
\item Profil
\item Avatar
\item Préférences
\item Langues
\item Notifications
\item Historique
\item Paramètres
\end{itemize}
\end{multicols}

\chapter{COMPANY API}

\section{Entités}
\begin{multicols}{2}
\begin{itemize}
\item Entreprise
\item Départements
\item Employés
\item Filiales
\item Agences
\item Logo
\item Identité graphique
\end{itemize}
\end{multicols}

\chapter{PROJECTS API}

\section{Fonctions}
\begin{multicols}{2}
\begin{itemize}
\item Création
\item Tâches
\item Planning
\item Livrables
\item Commentaires
\item Historique
\item Statistiques
\end{itemize}
\end{multicols}

\chapter{ORDERS API}

\section{Workflow}
\begin{enumerate}
\item Création
\item Validation
\item Production
\item Contrôle
\item Livraison
\item Archivage
\end{enumerate}

\chapter{MARKETPLACE API}

\section{Fonctions}
\begin{multicols}{2}
\begin{itemize}
\item Produits
\item Catégories
\item Stocks
\item Téléchargements
\item Licences
\item Commentaires
\item Évaluations
\end{itemize}
\end{multicols}

\chapter{ART GALLERY API}

\section{Entités}
\begin{multicols}{2}
\begin{itemize}
\item Tableaux
\item Illustrations
\item Affiches
\item Sculptures numériques
\item Collections
\item Éditions limitées
\item Certificats d'authenticité
\end{itemize}
\end{multicols}

\chapter{MEDIA API}

\section{Gestion}
\begin{multicols}{2}
\begin{itemize}
\item Images
\item Vidéos
\item PDF
\item Audio
\item Templates
\item Animations
\item Miniatures
\item Compression
\item Optimisation
\item Versioning
\end{itemize}
\end{multicols}

\chapter{AI CENTER API}

\section{Pilotage}
\begin{multicols}{2}
\begin{itemize}
\item Modèles IA
\item Agents
\item Prompts
\item Mémoire
\item Statistiques
\item Journaux
\end{itemize}
\end{multicols}

\chapter{AI AGENTS API}

\section{Exposition}
Chaque agent expose des services :
\begin{multicols}{2}
\begin{itemize}
\item Réception des tâches
\item Exécution
\item Compte-rendu
\item Journalisation
\item Transfert à un autre agent
\end{itemize}
\end{multicols}

\chapter{AUTOMATION API}

\section{Orchestration}
\begin{multicols}{2}
\begin{itemize}
\item Création automatique de projet
\item Attribution des tâches
\item Validation
\item Contrôle qualité
\item Livraison
\end{itemize}
\end{multicols}

\chapter{KNOWLEDGE BASE API}

\section{Fonctions}
\begin{multicols}{2}
\begin{itemize}
\item Import de documents
\item Indexation
\item Recherche
\item Gestion des versions
\item Suppression
\item Contrôle des accès
\end{itemize}
\end{multicols}

\chapter{RAG API}

\section{Fonctions}
\begin{itemize}
\item Interrogation des bases documentaires
\item Recherche des informations pertinentes
\item Alimentation des réponses des agents IA
\end{itemize}

\chapter{CHAT API}

\section{Conversations}
\begin{multicols}{2}
\begin{itemize}
\item Client ↔ IA
\item Client ↔ Support
\item Employé ↔ Employé
\item IA ↔ IA (journal interne)
\end{itemize}
\end{multicols}

\chapter{NOTIFICATION API}

\section{Support}
\begin{multicols}{2}
\begin{itemize}
\item Email
\item Notifications internes
\item SMS
\item Autres canaux configurés
\end{itemize}
\end{multicols}

\chapter{COMMUNICATION API}

\section{Canaux}
\begin{multicols}{2}
\begin{itemize}
\item WhatsApp Business
\item Gmail
\item Telegram
\item Facebook Messenger
\item Instagram Messaging
\item SMS
\item Signal
\item Discord
\item Slack
\item Microsoft Teams
\end{itemize}
\end{multicols}

\chapter{DELIVERY API}

\section{Options}
Le client choisit où recevoir ses livrables :
\begin{multicols}{2}
\begin{itemize}
\item Portail Client
\item Email
\item WhatsApp Business
\item Cloud privé
\item Lien sécurisé
\item Webhook
\end{itemize}
\end{multicols}

\chapter{PAYMENT API}

\section{Fonctions}
\begin{multicols}{2}
\begin{itemize}
\item Paiements
\item Abonnements
\item Remboursements
\item Facturation
\end{itemize}
\end{multicols}

\chapter{BILLING API}

\section{Types}
\begin{multicols}{2}
\begin{itemize}
\item Devis
\item Factures
\item Reçus
\item Avoirs
\item Abonnements
\end{itemize}
\end{multicols}

\chapter{AFFILIATE API}

\section{Gestion}
\begin{multicols}{2}
\begin{itemize}
\item Liens
\item QR Codes
\item Commissions
\item Retraits
\item Campagnes
\item Statistiques
\end{itemize}
\end{multicols}

\chapter{TRAINING API}

\section{Composants}
\begin{multicols}{2}
\begin{itemize}
\item Cours
\item Quiz
\item Certificats
\item Examens
\item Suivi de progression
\end{itemize}
\end{multicols}

\chapter{ANALYTICS API}

\section{Suivi}
\begin{multicols}{2}
\begin{itemize}
\item Trafic
\item Ventes
\item Campagnes
\item Utilisation des agents IA
\item Performances des modules
\end{itemize}
\end{multicols}

\chapter{SECURITY API}

\section{Mesures}
\begin{multicols}{2}
\begin{itemize}
\item Chiffrement
\item Gestion des secrets
\item Journalisation
\item Détection des anomalies
\item Rate limiting
\item Audit
\end{itemize}
\end{multicols}

\chapter{DEVELOPER CENTER}

\section{Espace}
\begin{multicols}{2}
\begin{itemize}
\item Documentation
\item Console de test
\item Génération de clés API
\item Suivi de consommation
\item Exemples d'intégration
\end{itemize}
\end{multicols}

\chapter{API DOCUMENTATION}

\section{Contenu}
\begin{itemize}
\item Description
\item Paramètres
\item Exemples
\item Codes d'erreur
\item Scénarios d'utilisation
\end{itemize}

\chapter{WEBHOOKS}

\section{Evénements}
\begin{multicols}{2}
\begin{itemize}
\item Commande créée
\item Projet terminé
\item Paiement validé
\item Livraison effectuée
\item Nouveau message
\end{itemize}
\end{multicols}

\chapter{SDK ENTERPRISE}

\section{Bibliothèques}
Prévoir des bibliothèques clientes pour faciliter l'intégration.

\chapter{VERSIONING}

\section{Principes}
Toutes les API doivent être versionnées :
\begin{itemize}
\item v1
\item v2
\end{itemize}
Les anciennes versions restent disponibles pour compatibilité.

\chapter{MONITORING}

\section{Indicateurs}
\begin{multicols}{2}
\begin{itemize}
\item Disponibilité
\item Temps de réponse
\item Erreurs
\item Taux d'utilisation
\item Consommation des ressources
\end{itemize}
\end{multicols}

\chapter{ROADMAP}

\section{Évolutivité}
L'architecture API doit permettre d'ajouter facilement de nouveaux services.

\chapter{AI CONNECTOR HUB}

\section{Concept}
Créer un Centre de Connecteurs IA capable de relier les agents à des services externes.

\section{Catégories}
\begin{multicols}{2}
\begin{itemize}
\item Messagerie
\item Stockage cloud
\item Suites bureautiques
\item CRM externes
\item ERP externes
\item Réseaux sociaux
\item Plateformes e-commerce
\end{itemize}
\end{multicols}

\section{Gestion}
Chaque connecteur gère :
\begin{itemize}
\item Authentification
\item Autorisations
\item Synchronisation
\item Erreurs
\item Journaux
\end{itemize}

\chapter{ORCHESTRATION DE LIVRAISON MULTICANAL}

\section{Moteur}
Développer un moteur de livraison capable de :
\begin{itemize}
\item Détecter les canaux autorisés par le client
\item Préparer automatiquement les livrables au bon format
\item Effectuer les envois
\item Suivre les confirmations de livraison
\item Gérer les échecs et les relances
\end{itemize}

\appendix

\chapter{Index}
\begin{enumerate}
\item Vision de l'Architecture
\item Architecture Microservices
\item API Gateway
\item Service Registry
\item Authentication API
\item Authorization API
\item Users API
\item Company API
\item Projects API
\item Orders API
\item Marketplace API
\item Art Gallery API
\item Media API
\item AI Center API
\item AI Agents API
\item Automation API
\item Knowledge Base API
\item RAG API
\item Chat API
\item Notification API
\item Communication API
\item Delivery API
\item Payment API
\item Billing API
\item Affiliate API
\item Training API
\item Analytics API
\item Security API
\item Developer Center
\item API Documentation
\item Webhooks
\item SDK Enterprise
\item Versioning
\item Monitoring
\item Roadmap
\item AI Connector Hub
\item Orchestration Livraison
\end{enumerate}

\end{document}
